%--------------------------------------------------------------------------------
%
%--------------------------------------------------------------------------------
\unnumberedChapter{UI Elements}

The Controller Dependent Code library abstracts the hardware platform, while the
LCS Runtime library abstracts communication between nodes. Yet another layer of
complexity remains in many embedded applications: the implementation of the user
interface.

Almost every controller project starts with a blinking LED and a push button.
Soon additional buttons are added, LEDs begin to blink, rotary encoders need to
be decoded and displays are required to present information to the user. Before
long, every application contains its own implementation of button debouncing,
long-press detection, display drivers and menu handling.

Although these requirements occur in many projects, they are frequently
implemented repeatedly from scratch. The resulting code is often tightly coupled
to the underlying hardware and difficult to reuse in subsequent projects.

The purpose of the \textbf{UI Elements} library is to provide reusable software
components for the most common user interface elements found in LCS hardware
modules. Instead of working with electrical signal levels or display controller
commands, applications interact with buttons, LEDs, rotary encoders and
character displays through a common, hardware independent interface.

The UI Elements library intentionally focuses on the individual building blocks
of a user interface. More sophisticated user interface concepts such as screen
management and menu navigation are built on top of these elements and are
described in the following chapter.

%--------------------------------------------------------------------------------
\unnumberedSection{Concept}

The UI Elements library follows the same design philosophy as the other LCS
libraries. Rather than exposing the underlying hardware, it presents the
application with higher level user interface objects.

For example, a push button is much more than a digital input pin. Mechanical
contacts bounce when they are pressed and released. Furthermore, applications
are rarely interested in the electrical signal itself. Instead, they want to
know whether the user performed a click, a double click or a long press.

Likewise, an LED is more than a digital output pin. Besides switching the LED on
and off, applications frequently require blinking patterns or status
indications. Similarly, rotary encoders should report rotational movement rather
than individual signal transitions, while displays should present text instead
of exposing controller specific communication protocols.

To achieve this abstraction, every UI element is implemented as an independent
software object. Internally, these objects use small state machines to analyze
their inputs and maintain their current state. The application periodically
calls a common \texttt{tick()} routine, which advances the state of all active
UI elements.

Whenever a relevant event is detected, such as a button click or a change of the
rotary encoder position, the corresponding callback function is invoked. The
application therefore reacts to user actions instead of continuously polling
hardware inputs.

Another important design goal is the complete separation of user interface
behavior from the underlying hardware implementation. Instead of directly
accessing controller pins, UI elements obtain their input values or update their
outputs through callback functions. This allows identical UI objects to operate
equally well with direct GPIO pins, shift registers, I/O expanders or any other
hardware interface capable of providing the required signals.

%--------------------------------------------------------------------------------
\unnumberedSection{Buttons}

Buttons are among the most frequently used user interface elements. Although
their hardware implementation is simple, their behavior is surprisingly
complex. Mechanical contacts bounce, users may press a button only briefly or
hold it down for an extended period, and many applications distinguish between
single clicks, double clicks and long presses.

The UIButton object encapsulates all these aspects in a compact state machine.
Each button instance maintains its own timing information and generates the
corresponding events automatically. The debounce interval as well as the timing
for click and long-press detection can be configured individually.

Applications do not poll the button continuously. Instead, callback functions
are registered for the desired events and are invoked whenever the corresponding
user action has been detected.

The following lines depict the button object. Note that only the public part of
is shown.

\lstset{language=c++, style=codesnippetstyle, numbers=none}
\begin{lstlisting}
struct UIButton : UIElements {

    public:

    UIButton( uint8_t rNum, bool activeLow = false );

    void      attachClick( UIButtonCbFunction functionId );
    void      attachDoubleClick( UIButtonCbFunction functionId );
    void      attachLongPressStart( UIButtonCbFunction functionId ) ;
    void      attachLongPressStop( UIButtonCbFunction functionId );
    void      attachDuringLongPress( UIButtonCbFunction functionId );
    void      attachGetDataFunction( UIGetDataFunction functionId );

    void      setActiveLow( bool val );
    bool      isLongPressed( );
    uint32_t  getDurationPressedMillis( );
    uint32_t  getPressedMillisSinceStart( );
    uint8_t   getResId( );

    void      reset( );
    void      processTick( );
};
\end{lstlisting}
\FloatBarrier

The actual retrieval of the button state is itself abstracted by means of a
callback function. While this may initially appear unnecessarily flexible, it
allows the same button object to operate with many different hardware
implementations.

A button may be connected directly to a controller GPIO pin, be part of a serial
shift register or reside on an external I/O expander such as the MCP23017. From
the point of view of the button state machine, all these hardware
implementations appear identical.

The button identifier therefore does not necessarily correspond to a physical
controller pin. It merely represents an identifier understood by the callback
function responsible for obtaining the current button state.

%--------------------------------------------------------------------------------
\unnumberedSection{Rotary Encoders}

Rotary encoders extend the button concept by reporting rotational movement
instead of discrete button events. A typical incremental encoder provides two
quadrature signals whose phase relationship determines the direction of
rotation.

The UIEncoder object continuously evaluates these signals and reports clockwise
or counter-clockwise movement through callback functions. Just like buttons, the
actual signal acquisition is delegated to a hardware callback, allowing rotary
encoders connected through GPIO pins, shift registers or I/O expanders to be
handled identically.

\lstset{language=c++, style=codesnippetstyle, numbers=none}
\begin{lstlisting}
struct UIEncoder : UIElements {

    public:

    UIEncoder(  uint8_t rNum,
                int     lower     = INT_MIN, 
                int     upper     = INT_MAX, 
                bool    activeLow = false );

    void        reset( );
    void        processTick( void );

    void        setLimits( int lower, int upper );
    int         getLowerLimit( );
    int         getUpperLimit( );
    int         getPosition( );
    uint8_t     getResId( );

    void        setPosition( int  newPosition,
                             bool suppressCallback = false );
                                             
    uint32_t    getMillisBetweenRotations( );
    void        attachPositionChanged( UIEncoderCbFunction functionId );
    void        attachGetDataFunction( UIGetDataFunctionPair functionId );
};
\end{lstlisting}
\FloatBarrier

%--------------------------------------------------------------------------------
\unnumberedSection{LEDs}

LEDs provide visual feedback to the user and therefore complement buttons and
rotary encoders. The UILED object represents an indicator rather than a simple
digital output.
Besides switching the LED on and off, the object supports common blinking
patterns controlled by an internal state machine. Applications therefore specify
the desired behavior instead of periodically toggling the output pin
themselves.

\lstset{language=c++, style=codesnippetstyle, numbers=none}
\begin{lstlisting}
struct UILed : UIElements {

    public:

    UILed( uint8_t rNum );

    void processTick( );
    void attachSetDataFunction( UISetDataFunction functionId );

    bool isOn( );
    bool isOff( );
    void setOn( );
    void setOff( );
    void setVal( bool val );
    void toggle( );
    void blink( );

    static void setBlinkIntervalMillis( uint32_t val );
};
\end{lstlisting}
\FloatBarrier

As with all other UI elements, updating the actual hardware output is delegated
to a callback function, making the implementation independent of the underlying
hardware interface.

%--------------------------------------------------------------------------------
\unnumberedSection{Displays}

Character displays exist in many different forms. Traditional LCD modules,
graphical OLED displays and other display technologies all differ in their
electrical interface and controller capabilities. Nevertheless, many embedded
applications require only a comparatively small common feature set.

The UIDisplay object therefore models a display as a rectangular matrix of ASCII
characters. Applications position a cursor, write characters and clear the
display without any knowledge of the underlying display technology.

Controller specific capabilities such as adjustable brightness, contrast,
backlight control or multiple font sizes remain available through the respective
display driver. The common interface intentionally covers only the operations
shared by all supported display types.

\lstset{language=c++, style=codesnippetstyle, numbers=none}
\begin{lstlisting}
struct UIDisplayLcdI2C : public UIDisplay {

    public:

    UIDisplayLcdI2C( uint8_t dType, 
                     uint8_t rNum, 
                     uint8_t i2cAdr = 0x27 );

    void    displayOn( );
    void    displayOff( );
    void    setCursor( uint8_t col, uint8_t row );
    uint8_t print( const char *s );
    uint8_t print( char ch );
    void    clear( );
    void    clearLine( uint8_t row );

    private: 

    LcsLcdDisplay *lcd = nullptr;
};\end{lstlisting}
\FloatBarrier

The display class shown is actually the base class. On top there are inheriting 
classes for an LCS display and an OLED display. 

%--------------------------------------------------------------------------------
\unnumberedSection{Summary}

The UI Elements library provides a collection of reusable building blocks for
implementing user interfaces on LCS hardware modules. Rather than exposing
controller pins or display specific communication protocols, the library
presents the firmware with objects representing the common user interface
elements found in embedded systems.

Buttons, LEDs, rotary encoders and displays each encapsulate both the underlying
hardware access and the behavioral aspects associated with these devices.
Mechanical button debouncing, click detection, blinking patterns and encoder
decoding are all handled internally by dedicated state machines. Applications
therefore react to meaningful user events instead of low-level electrical signal
changes.

A fundamental design goal is the complete separation of user interface behaviour
from the physical hardware implementation. All interaction with the underlying
controller or external interface hardware is performed through callback
functions. Consequently, identical UI objects can operate with direct GPIO pins,
shift registers, I/O expanders or other hardware interfaces without requiring
any modifications.

The UI Elements library concentrates on the individual building
blocks of a user interface. More advanced concepts such as screen management,
menu navigation and complete operator interfaces are constructed from these
elements. The next chapter introduces this additional abstraction layer and
shows how complete user interfaces can be assembled from the UI elements
presented here.
